% Source-aligned English for printed p.147.
\SourcePage{147}

\LayerTarget{tohoku:construction:hyperhomology-resolution:a:p2}
\noindent an injective object.

\LayerTarget{tohoku:construction:hyperhomology-resolution:b}
b) Take an ``injective resolution'' of $K$ in the sense of [6, Chapter XVII].
The quotation marks are necessary because the terminology of [6] is plainly
ambiguous. Thus the $L^{p,q}$ are injective and, for fixed $p$, taking the
cycles, boundaries, or cohomology of $L^{p,*}$ with respect to the first
differential gives an injective resolution of the object formed by the cycles,
boundaries, or cohomology, respectively, in $K^p$.

\LayerTarget{tohoku:construction:hyperhomology-identification}
If $L$ is as above, $H^*F(L)$ is essentially independent of the choice of $L$
and identifies with $R^*(F\mathbin{\circ}H^0)(K)$. Here $H^0$ is regarded as a
covariant left exact functor $\mathfrak K\longrightarrow C$, and
$F\mathbin{\circ}H^0$ is the composite $\mathfrak K\longrightarrow C'$. To see
this, take an injective resolution $L'$ of $K$ in the sense of a). There is a
homomorphism, unique up to homotopy, from the resolution $L$ to $L'$, hence a
homomorphism $F(L)\longrightarrow F(L')$ inducing an isomorphism
\[
{}_{\mathrm I}E_2^{p,q}(F(L))\longrightarrow
{}_{\mathrm I}E_2^{p,q}(F(L')).
\]
Both sides identify with $H^p(R^qF(K))$, so this induces an isomorphism from
$HF(L)$ to $HF(L')$. The latter is described by the second spectral sequence
of the bicomplex $F(L')$: one has $H_{\mathrm I}^q(F(L'))=0$ for $q>0$, and
hence
\[
H^nF(L')=H^n(FH^0(L'^{*,n})).
\]
The right side is, by definition, $R^n(F\mathbin{\circ}H^0)(K)$. This defines
and computes the \textit{hyperhomology}
$\mathfrak R F(K)=\mathfrak R^*(FH^0)(K)$ of $F$ with respect to $K$, together
with the first spectral sequence, whose initial term is
${}_{\mathrm I}E_2^{p,q}F(K)=H^p(R^qF(K))$ and which abuts to it.

If $L$ is instead an ``injective resolution'' as in b), the second spectral
sequence of $F(L)$ is essentially independent of $L$, since $L$ is unique up
to homotopy equivalence. It abuts to $\mathfrak R F(K)$ and has the initial
term in (2.4.2). It would in fact suffice to require
$R^nF(L^{p,q})=0$ for $n>0$, rather than injectivity of $L^{p,q}$, for the
second spectral sequence of $F(L)$ to be the one in (2.4.2).

\LayerTarget{tohoku:remark:derived-spectral-functor-scope}
If $K$ is in nonnegative degrees, both derived spectral sequences of $F$ are
cohomological. They can also be defined without assuming that $K$ is bounded
below, provided $F$ has finite injective dimension, that is, $R^pF=0$ for a
suitable $p$. This does not seem to be contained in [6]; since it will not be
used below, we merely record it without proof. One can of course define left
derived spectral functors of $F$ when $C$ has enough projective objects, and
one can also treat contravariant functors. Reference [6] treats a
multifunctor, but we shall not need that case.

\LayerTarget{tohoku:construction:composite-functor-spectral-sequences:p1}
Let $C,C',C''$ be abelian categories and let
$F:C\longrightarrow C'$ and $G:C'\longrightarrow C''$ be covariant functors.
Suppose that every object of $C$ or $C'$ is isomorphic to a subobject of an
injective object, so that the right derived functors of $F,G$, and $GF$ may be
formed. We seek relations among them. Let $A\in C$ and let $C(A)$ be the
complex associated to a right resolution of $A$ by injective objects. The
complex $F(C(A))$ in $C'$ is determined up to homotopy as $C(A)$ varies.

\TohokuContinuation{148}
